\section*{الفصل \ref{ch:b3:microbiomes} --- الميكروبيومات والتعايشات}

\begin{solution}{exo:b3:microbiomes:1}
\omterm{def:b3:microbiomes:symbiosis}{التقايض}: يربح الطرفان --- البقولية والريزوبيوم، والمرجان
والطحلب. \omterm{def:b3:microbiomes:symbiosis}{والمعايشة}: يربح أحدهما ولا يتأثر الآخر --- ومنه
معظم بكتيريا الأمعاء التي تقتات ما يعجز المضيف عن هضمه.
والتطفّل: يربح أحدهما على حساب الآخر --- ومنه
\emph{Wolbachia} تقتل الأجنة الذكور. وقد ينتقل الزوج نفسه:
فمعايش معوي يعبر الجدار يصير ممرضًا، وريزوبيوم يكف عن
التثبيت يصير طفيليًا على العقدة، وطحلب تحت الإجهاد الحراري
يضر المرجان الذي آواه.
\end{solution}

\begin{solution}{exo:b3:microbiomes:2}
$H = -(0.5\ln 0.5 + 0.3\ln 0.3 + 0.1\ln 0.1 + 2\times 0.05\ln 0.05) =
0.347 + 0.361 + 0.230 + 0.300 = 1.24$؛ و$e^{H} = 3.4$ نوعًا
فعالًا. و$D = 0.25 + 0.09 + 0.01 + 0.0025 + 0.0025 = 0.355$؛
و$1/D =
2.8$.
\end{solution}

\begin{solution}{exo:b3:microbiomes:3}
يخبر مسح 16S بمن هو موجود، إلى حدود الجنس، وبعدد قراءاته
النسبي؛ ولا يقول شيئًا عن المورّثات التي يحملها، ويفوته
الفطور والفيروسات (بادئات أخرى)، ويتحيزه عدد النسخ
والبادئات. وأما ميتاجينوميات البندقية فتخبر بالمورّثات
والوظائف، وبالجينومات الكاملة والسلالات إذا بلغ العمق
كفايته؛ وهي أغلى في العينة الواحدة، ويغرقها DNA المضيف في
عينات النسج، ولا تزال عاجزة عن قول أي المورّثات معبَّر عنها
وأي الخلايا حية.
\end{solution}

\begin{solution}{exo:b3:microbiomes:4}
تخمير الألياف إلى أحماض دهنية قصيرة السلسلة تغذي القولون
والمضيف؛ وتركيب فيتامين~K وفيتامينات B؛ وتحويل الأحماض
الصفراوية والأدوية؛ ومقاومة استعمار الممرضات؛ وتدريب
الجهاز المناعي. والاختلال: التهاب قولون
\emph{Clostridioides difficile} بعد الصادات (ومما اتُّهم به
أيضًا: داء الأمعاء الالتهابي، والبدانة).
\end{solution}

\begin{solution}{exo:b3:microbiomes:5}
$\binom{20}{5} = \num{15504}$. والنوع ذو $10$: $1 - \binom{10}{5}/
\num{15504} = 1 - 252/\num{15504} = 0.984$؛ وذو $6$: $1 - 2002/\num{15504}
= 0.871$؛ وذو $3$: $1 - 6188/\num{15504} = 0.601$؛ وذو $1$: $1 -
\num{11628}/\num{15504} = 0.25$. و$E[S_{5}] = 2.7$ نوعًا. والتغطية:
نوع مفرد واحد، $1 - 1/20 = 0.95$.
\end{solution}

\begin{solution}{exo:b3:microbiomes:6}
$400\times 10^{11} = 4\times 10^{13}$ بكتيريا، نحو \qty{40}{g}.
والمورّثات: $1000$ نوع $\times$ \num{4000}، ناقصًا التداخل ---
أي من رتبة $3\times 10^{6}$ مورّثة متمايزة، أكثر من مئة ضعف
مورّثات الإنسان \num{20000}.
\end{solution}

\begin{solution}{exo:b3:microbiomes:7}
يهدم الأكسجين \omterm{def:b3:microbiomes:nodule}{النتروجيناز}، غير أن البكتيرويدات مضطرة إلى
التنفس لتصنع ستة عشر ATP لكل N$_{2}$. وتشكّل قشرة العقدة
حاجز انتشار يحد من دخول الأكسجين، ويربط \omterm{def:b3:microbiomes:nodule}{الليغهيموغلوبين}،
عند تركيز ميلي مولاري، الأكسجينَ الداخل فيبقى التركيز الحر
قرب \qty{10}{nM} مع بقاء الدفق إلى البكتيرويدات عاليًا.
واللون الوردي هو أكسي-ليغهيموغلوبين؛ وأما العقدة الخضراء
فقد تحلل ليغهيموغلوبينها --- فهي هرمة لم تعد تثبّت.
\end{solution}

\begin{solution}{exo:b3:microbiomes:8}
يفشل الاثنان: فلا عقد ولا \omterm{def:b3:microbiomes:mycorrhiza}{ميكوريزا شجيرية}، لأن ذروة
الكالسيوم هي الإشارة المشتركة في \omterm{def:b3:microbiomes:mycorrhiza}{سبيل التعايش المشترك}.
فالسبيل إذن أقدم من تكوين العقد --- إذ بُني \omterm{def:b3:microbiomes:symbiosis}{للتعايش} الفطري
قبل أربعمئة مليون سنة ثم جنّدته البقوليات \omterm{def:b3:microbiomes:nodule}{للريزوبيا}
لاحقًا.
\end{solution}

\begin{solution}{exo:b3:microbiomes:9}
عرّض كيرز بعض عقد فول الصويا لجو من الأرغون والأكسجين
يتعذر فيه على \omterm{def:b3:microbiomes:nodule}{الريزوبيا} تثبيت النتروجين، بينما ثبّتت بقية
عقد النبات تثبيتًا سويًا؛ فقطع النبات إمداد الأكسجين عن
العقد غير المثبّتة فتكاثرت ريزوبياها نحو النصف. ولولا هذا
التمييز لتكاثرت \omterm{def:b3:microbiomes:nodule}{الريزوبيا} التي تتخطى التثبيت المكلف أسرع
من المثبّتة، ولانتشرت في الجمهرة، ولانتهت النباتات إلى
إيواء بكتيريا لا تعطي شيئًا: فينهار \omterm{def:b3:microbiomes:symbiosis}{التقايض}.
\end{solution}

\begin{solution}{exo:b3:microbiomes:10}
$\binom{N-N_{i}}{n}/\binom{N}{n} = \prod_{j=0}^{n-1}(N - N_{i} - j)/(N -
j)$، وكل عامل فيه دون $1$؛ والانتقال من $n$ إلى $n + 1$ يضرب
بعامل آخر دون $1$، فيهبط احتمال الإخفاق ويرتفع احتمال
الحضور مع $n$؛ وعند $n = N$ يكون البسط $\binom{N -
N_{i}}{N}$ صفرًا فتحضر الأنواع كلها، فنحصل على $S$. وعند $N_{i},
n \ll N$ يقارب كل عامل $1 - N_{i}/N$، ويقارب الجداء $(1 -
N_{i}/N)^{n} \approx e^{-nN_{i}/N} \approx (1 - n/N)^{N_{i}}$ ---
وهو تقريب المعاينة مع الإرجاع.
\end{solution}

\begin{solution}{exo:b3:microbiomes:11}
الذكر طريق مسدود لمتعايش لا يسافر إلا في البيوض، فأي صفة
تحوّل الذكور إناثًا، أو تقتلها لمصلحة أخواتها، أو تعوق
الإناث غير المصابة، ترفع انتقال المتعايش. وأما عدم التوافق
السيتوبلازمي: فنطاف الذكور المصابة تُعدَّل بحيث تموت
البيضة الملقحة ما لم تحمل البيضة \emph{Wolbachia} نفسها،
فتنقذها. ومن ثم تخسر الإناث غير المصابة الخلف الذي تحمله
من ذكور مصابة، ولا تخسر الإناث المصابة شيئًا؛ وتنمو
الأفضلية مع تواتر الخمج، فينتشر الخمج فوق عتبة معينة حتى
التثبيت ولو وضعت الإناث المصابة بيوضًا أقل قليلًا --- إذ
ترجح الكلفةَ نسبةُ حضنات الإناث غير المصابة التي تُتلف.
\end{solution}

\begin{solution}{exo:b3:microbiomes:12}
يمر المتعايش الداخلي بعنق زجاجة من خلايا قليلة في كل جيل
من المضيف ولا يتأشب قط مع خارجه: فيثبّت الانجراف طفرات
ضارة ضررًا طفيفًا، ومنها الحذوف، ولا مصدر لمورّثات بديلة
(سقّاطة مولر). والمورّثات التي يوفر المضيف منتجاتها، أو
التي تخدم حياة حرة --- \omterm{def:b3:virology:virus}{الغلاف}، والحركة، والتنظيم، ومعظم
الإصلاح --- لا تُفتقد فتذهب أولًا؛ وأما آلة الترجمة
والمورّثات التي تصنع ما يحتاج إليه المضيف (الأحماض
الأمينية الأساسية في \emph{Buchnera}) فتبقى أطول. وأقاربه
الحرة جمهراتها ضخمة، وفيها تأشب وتدفق مورّثات، والانتخاب
يصون جينوماتها. وعكس الاتجاه يحتاج إلى مورّثات جديدة من
الخارج --- وهو متعذر في العزلة --- فيستبدل المضيف بدل ذلك
متعايشًا متآكلًا بآخر جديد، كما فعلت عدة سلالات حشرية.
\end{solution}

\begin{solution}{pb:b3:microbiomes:1}
\textbf{1.} $400\times 10^{11} = 4\times 10^{13}$ بكتيريا، \qty{40}{g}.
\textbf{2.} جينوم من \num{4000} مورّثة يقارب \qty{4}{Mb}، أي \qty{4.4}{fg}:
$4\times 10^{13}\times 4.4\times 10^{-15} = \qty{0.18}{g}$ من DNA
البكتيري، في مقابل $3\times 10^{13}\times 6.4\times 10^{-12} = \qty{0.19}{g}$
من DNA البشري --- أي متساويان تقريبًا.
\textbf{3.} $1000\times 0.7\times 4000 + 0.3\times 4000 \approx 2.8\times
10^{6}$ مورّثة متمايزة، أي نحو $140$ ضعف مورّثات جينوم الإنسان.
\textbf{4.} $30\times 10 = \qty{300}{mmol}$؛ و$\times 0.9 = \qty{270}{kJ}$؛
أي \qty{3}{\%} من الحمية.
\textbf{5.} لا --- إذ لا تفسر \qty{3}{\%} نسبة \qty{30}{\%}. والفئران
الخالية من الجراثيم تمتص أيضًا امتصاصًا أقل كفاءة،
وهرموناتها المعوية وتخزينها للشحم وتوليدها للحرارة متبدلة،
وشهيتها مختلفة: \omterm{def:b3:microbiomes:symbiosis}{فالميكروبيوم} يعمل في فيزيولوجيا المضيف، لا
مورّدًا للوقود فحسب.
\textbf{6.} انقسام يوازنه الطرح: جمهرة مستقرة، تُستبدل كل
خلية منها يوميًا. وبعد قتل \qty{99.9}{\%} يبقى $4\times 10^{10}$
وتعيد عشر مضاعفات الأعداد في نحو عشرة أيام --- غير أن
الناجين والأسرع نموًا يهيمنون، ويتبدل التركيب (خلل توازن)،
وفي الفجوة يمكن لغازٍ مثل \emph{C.
difficile} أن يستقر.
\textbf{7.} تتقاسم الأنواع الثانوية $0.4$ على $177$: أي $p = 0.00226$
لكل منها. و$H
= -(0.3\ln 0.3 + 0.2\ln 0.2 + 0.1\ln 0.1 + 0.4\ln 0.00226) = 0.361 +
0.322 + 0.230 + 2.437 = 3.35$؛ و$e^{H} \approx 28$.
\textbf{8.} $D = 0.09 + 0.04 + 0.01 + 177\times (0.00226)^{2} = 0.141$؛
و$1/D = 7.1$. \omterm{def:b3:microbiomes:diversity}{ودليل سيمبسون} يربّع الوفرات، فتقرره الأنواع
الثلاثة المهيمنة ولا تكاد النادرة تُحسب؛ وأما شانون فيعطي
الأنواع النادرة وزنًا أكبر.
\textbf{9.} $1 - 60/2000 = 0.97$: أي نحو $30$ قراءة في الألف
تنتمي إلى أنواع لم تُرَ بعد.
\textbf{10.} يرتفع الغنى مع حجم العينة ما دامت الأنواع النادرة
تظهر؛ إذ تبلغ \num{10000} قراءة أنواعًا تفوتها \num{2000}. رقّق
العينة الكبيرة إلى \num{2000} قراءة بالمبرهنة (أو بالمعاينة
الجزئية) --- وينبغي أن تعطي نحو $180$ إن كانت الجماعة
نفسها --- ثم قارن عند العمق نفسه.
\textbf{11.} عند $N_{i} = 1$: $1 - \binom{1999}{1000}/\binom{2000}{1000} =
1 - (2000 - 1000)/2000 = 0.5$. وعند $N_{i} = 5$: نحو $1 - 0.5^{5} = 0.97$.
\textbf{12.} $e^{2.0} = 7.4$ نوعًا فعالًا؛ و$D \ge 0.5^{2} = 0.25$،
أي نحو $0.27$: جماعة فيها نوع واحد مهيمن وبضع عشرات من
الأنواع الثانوية --- وهي الأمعاء المنهارة القليلة التنوع
التي يستقر فيها \emph{C.~difficile}.
\textbf{13.} $150\times 8 = \qty{1200}{kg}$ من السكر في الهكتار،
أي \qty{10}{\%} من \qty{12}{t} المركَّبة ضوئيًا.
\textbf{14.} $150\,000/28 = 5360$ مولًا من N$_{2}$؛ و$16\times 5360 =
\num{85700}$ مولًا من ATP.
\textbf{15.} $\num{85700}/30 = 2860$ مولًا من الغلوكوز، أي \qty{514}{kg} ---
نحو \qty{43}{\%} من السكر المدفوع. والباقي يبني العقد
والبكتيرويدات ويصونها، ويوفر الإلكترونات الثمانية لكل
N$_{2}$ مُرجِعًا، ويمثّل النشادر.
\textbf{16.} $1200\times 16 = \qty{19200}{MJ}$ مقابل \qty{150}{kg} من N:
أي \qty{128}{MJ/kg}، نحو أربعة أضعاف \qty{35}{MJ/kg} في المصنع ---
غير أنها مدفوعة بضوء الشمس، في الحقل، بلا وقود.
\textbf{17.} المثبّتات $0.9\times 1 = 0.9$؛ والغشّاشات $0.1\times 0.5 =
0.05$؛ فنسبة الغشّاشات $0.05/0.95 = 5.3\,\%$ --- أي تنصّفت في موسم واحد.
\textbf{18.} الغشّاشات $0.1\times 1.2^{10} = 0.62$ في مقابل مثبّتات $0.9$:
أي \qty{41}{\%} بعد عشرة مواسم، وهي مرتفعة نحو الهيمنة.
فالعقوبات تقلب الانتخاب على الغشّاش؛ ومن دونها يتآكل \omterm{def:b3:microbiomes:symbiosis}{التقايض}.
\textbf{19.} $10^{6}\times \qty{20}{pg} = \qty{20}{\micro g}$ من الكربون
في cm$^{2}$ يوميًا؛ ويبلغ المرجان \qty{18}{\micro g}.
\textbf{20.} التنفس \qty{15}{\micro g}: أي فائض \qty{3}{\micro g}
في cm$^{2}$ يوميًا للنمو، وللقالب العضوي في التكلس،
وللمخاط والأعراس.
\textbf{21.} عند \qty{5}{\%} من الطحالب، \qty{0.9}{\micro g} يوميًا
في مقابل \qty{15}{\micro g} متنفَّسة: أي عجز \qty{14}{\micro g} في
اليوم؛ ومخزون \qty{300}{\micro g} يدوم نحو ثلاثة أسابيع.
\textbf{22.} أربعة أسابيع عند $+2$ أسوأ: فالضرر اللاحق بأجهزة
الطحالب الضوئية يرتفع ارتفاعًا حادًا، لا خطيًا، مع الحرارة.
ومع ذلك يصلح المقياس إنذارًا لأن الابيضاض يكامل الإجهاد
التراكمي، ولأن الفرق بين المسارين أصغر من الارتياب في
العتبة؛ فهو دليل تجريبي، معاير على ابيضاض مرصود.
\textbf{23.} يمكن لمرجان مجهَد أن يميل إلى أنواع طحلبية أكثر
تحملًا للحرارة موجودة أصلًا بوفرة منخفضة، أو أن يكتسبها من
الماء بعد الابيضاض؛ والطحالب المتحملة تعطي كربونًا أقل،
فينمو المرجان أبطأ، والتحمل المكتسب درجة أو درجتان --- أقل
من الاحترار المتوقع --- بينما تفوق وتيرة الاحترار تأقلم
المرجان نفسه.
\textbf{24.} $10^{10}$ cm$^{2}$ $\times$ \qty{20}{\micro g} $=
\qty{200}{kg}$ من الكربون يوميًا، أي نحو \qty{73}{t} سنويًا.
\textbf{25.} نحو \qty{3}{\%} من الحمية من تخمير \omterm{def:b3:microbiomes:symbiosis}{الميكروبيوم}؛
و\qty{97}{\%} تغطية لعينة التسلسل؛ و\qty{10}{\%} من التركيب
الضوئي دُفعت مقابل نتروجين محصول الفول.
\end{solution}
